\section{Jefferson and Wilkinson: Access, Usefulness, Protection}

The surviving record suggests that Jefferson and Wilkinson were neither strangers nor intimate friends. They occupied a more typical but harder-to-name early republican relationship: proximity built from usefulness, then tested by crisis.

The early correspondence shows Wilkinson presenting himself as a man worth keeping in Jefferson's circle because he could provide western knowledge, objects, and personnel. In 1802 he sent river-water samples and fossils because they ``may not be unacceptable'' to Jefferson \citep{wilkinson1802}. In 1805 he forwarded a buffalo-hide map and vouched for Captain Amos Stoddard as a man of intelligence, closing with ``perfect respect'' and ``sincere attachment'' \citep{wilkinson1805}. Even the 1804 dinner note is revealing: social access existed, but it was immediately absorbed back into public business through Humboldt, Texas, and Mexico \citep{wilkinson1804jefferson}. The relationship was personal enough for dinner, but the dinner itself seems to have been another venue for utility.

That pattern became explicit during the Burr crisis. Jefferson's letter of 3 February 1807 is the clearest proof that Wilkinson's survival was tied to presidential need. Jefferson reassured him that ``we ... have not failed to strengthen the public confidence in you'' and promised that he would be ``cordially supported in the line of your duties'' \citep{jefferson1807}. That is not the language of detached neutrality. It is the language of a president who needs his general's credibility to hold while Burr's conspiracy is being narrated to the nation.

Wilkinson understood the asymmetry perfectly. His answer of 20 March 1807 thanked Jefferson for the ``sweet Solace'' of sympathy and described presidential support as ``the warrant of part confidence \& the pledge of future trust'' \citep{wilkinson1807a}. In April he went further and looked ``for Orders \& direction'' concerning ``time \& place'' \citep{wilkinson1807b}. Wilkinson was not writing as a friend seeking consolation. He was writing as a compromised subordinate who needed executive endorsement to remain a usable instrument.

The 1811 exchange sharpens the point. Wilkinson's memoir letter to Jefferson is drenched in persecution and honor. He writes under ``the pressure of my persecutions'' and presents his memoirs as a weapon against public falsehood \citep{wilkinson1811}. Jefferson's answer is warm but carefully framed. He dismisses the ``under my thumb'' slander, says he believes Wilkinson because of his own knowledge of Burr and Wilkinson's conduct, and hopes the general's ``Argosy may ride out'' the storm \citep{jefferson1811}. The letter matters because Jefferson refuses to join the chorus of suspicion. Yet even here the belief is grounded in political history and conduct, not in intimate testimony about character.

\begin{table}[t]
\centering
\small
\begin{tabular}{lrrrr}
\toprule
Author & Warmth & Utility & Defense & Total \\
\midrule
James Wilkinson & 4 & 5 & 6 & 15 \\
Thomas Jefferson & 2 & 4 & 2 & 8 \\
\bottomrule
\end{tabular}
\caption{Lexicon counts from excerpted correspondence and summaries. Wilkinson's surviving voice is especially saturated with defense language, while Jefferson's is warmer only when folded into judgment and utility.}
\label{tab:themes}
\end{table}


The coded lexicon makes the same point in simpler terms. Wilkinson's excerpted voice in this corpus carries more defense terms than warmth terms, while Jefferson's language balances warmth with utility and judgment. That asymmetry fits the documentary pattern. Wilkinson needed Jefferson more urgently than Jefferson needed Wilkinson personally; Jefferson needed Wilkinson primarily when western policy, Burr, or later retrospective memory made him useful. The 1818 correspondence confirms the shift. Jefferson's estimate of Wilkinson's services remained ``unaltered,'' and Wilkinson treasured the ``approbation \& esteem'' as recognition long denied by the public \citep{jefferson1818,wilkinson1818}. What remained, then, was not intimacy in the ordinary sense. It was a durable politics of remembered service.
